\documentclass[12pt, a4paper]{article}

% ─── Paketler ───────────────────────────────────────────────────────────────
\usepackage[utf8]{inputenc}
\usepackage[T1]{fontenc}
\usepackage[turkish]{babel}
\usepackage{geometry}
\usepackage{graphicx}
\usepackage{amsmath}
\usepackage{listings}
\usepackage{xcolor}
\usepackage{hyperref}
\usepackage{booktabs}
\usepackage{float}
\usepackage{caption}
\usepackage{subcaption}
\usepackage{fancyhdr}
\usepackage{titlesec}
\usepackage{enumitem}
\usepackage{parskip}
\usepackage[export]{adjustbox}


\geometry{
  top=2.5cm, bottom=2.5cm,
  left=3cm, right=2.5cm
}

% ─── Kod renk şeması ────────────────────────────────────────────────────────
\definecolor{codegreen}{rgb}{0,0.6,0}
\definecolor{codegray}{rgb}{0.5,0.5,0.5}
\definecolor{codepurple}{rgb}{0.58,0,0.82}
\definecolor{backcolour}{rgb}{0.95,0.95,0.92}

\lstdefinestyle{cstyle}{
  backgroundcolor=\color{backcolour},
  commentstyle=\color{codegreen},
  keywordstyle=\color{blue},
  numberstyle=\tiny\color{codegray},
  stringstyle=\color{codepurple},
  basicstyle=\ttfamily\footnotesize,
  breakatwhitespace=false,
  breaklines=true,
  captionpos=b,
  keepspaces=true,
  numbers=left,
  numbersep=5pt,
  showspaces=false,
  showstringspaces=false,
  showtabs=false,
  tabsize=2,
  language=C
}

% ─── Üst/alt bilgi ──────────────────────────────────────────────────────────
\pagestyle{fancy}
\fancyhf{}
\rhead{INF 340 Mikroişlemciler}
\lhead{Akıllı Ekosistem -- Son Rapor}
\rfoot{\thepage}

% ─── Başlık ayarları ────────────────────────────────────────────────────────
\titleformat{\section}{\large\bfseries}{}{0em}{\thesection\quad}
\titleformat{\subsection}{\normalsize\bfseries}{}{0em}{\thesubsection\quad}

% ────────────────────────────────────────────────────────────────────────────
\begin{document}

% ─── Kapak ──────────────────────────────────────────────────────────────────
\begin{titlepage}
  \centering
  \vspace*{1.5cm}
  {\Large \textbf{INF 340 MİKROİŞLEMCİLER}\par}
  \vspace{1.5cm}
  {\LARGE \textbf{Akıllı İçecek Otomasyonu}\par}
    \vspace{0.5cm}
  {\large Unutacağınız bir akıllı ev sistemi değil, bir DENEYİM!\par}
  \vspace{2cm}
  \begin{tabular}{ll}
    \textbf{Resul Sanberk Kalko}  & 23401068 & Soğuk Algılama ve Müzik Açma\\[4pt]
    \textbf{Ayşe Selen Göğüş} & 23401063 & Temassız Jest Kontrolü ile Müzik Değiştirme\\[4pt]
    \textbf{Erdem Ertan}         & 22401986 & Akustik Lamba Kontrolü\\
  \end{tabular}
  \vfill
  {\large Mayıs 2026\par}
\end{titlepage}

% ─── İçindekiler ────────────────────────────────────────────────────────────
\tableofcontents
\newpage

% ════════════════════════════════════════════════════════════════════════════
\section{Proje Özeti}
% ════════════════════════════════════════════════════════════════════════════

Bu proje kapsamında, mikrodenetleyici tabanlı bir akıllı ekosistem sistemi
tasarlanmıştır. Sistem; termodinamik ve ışık sensörleri bir araya
getirerek üç temel işlev üzerinden yönetebilen bir
sisteme dönüştürmektedir:

\begin{enumerate}[leftmargin=*, label=\arabic*.]
  \item \textbf{Soğuk Algılama ve Müzik Açma:} İçeceğin sıcaklığı hedef değere
        ulaştığında Bluetooth üzerinden otomatik müzik başlatılması.
  \item \textbf{Temassız Jest Kontrolü:} El jestleriyle çalan şarkının
        ileri/geri değiştirilmesi.
  \item \textbf{Akustik Lamba Kontrolü:} El çırpma sesiyle ortam lambasının açılıp
        kapatılması ve parlaklığının ayarlanması.
\end{enumerate}

% ════════════════════════════════════════════════════════════════════════════
\section{Sistem Mimarisi}
% ════════════════════════════════════════════════════════════════════════════

ESP32 ve Arduino Uno mikrodenetleyici bazlı mimari,  Spotify API desteğiyle ortak noktada buluşan birimler en sonunda istenilen deneyimi oluşturmuştur.\\ \\

Sıcaklık ölçen sensörler sayesinde içecek optimal sıcaklığa ulaştığında bilgilendirme amaçlı Bluetooth bağlantısı ile açılan müzik ışık sensörü düzeneğini kullanıma hazır hale getirir. Işık sensörleri ile ortam parlaklık seviyesi ölçülerek jest anı algılanıp olası ışık değişimlerine karşın bu parlaklık eşik değeri her 30 saniyede tekrar ölçülmektedir. Her iki sistemdeki durum bilgilendirmeleri aynı zamanda OLED ekranlarla görünülebilir kılınmıştır.


% ════════════════════════════════════════════════════════════════════════════
\section{Kullanılan Donanım}
% ════════════════════════════════════════════════════════════════════════════

\begin{table}[H]
  \centering
  \caption{Donanım Bileşenleri}
  \label{tab:donanim}
  \begin{tabular}{lll}
    \toprule
    \textbf{Bileşen}          & \textbf{Model}      & \textbf{Arayüz} \\
    \midrule
    Ana Kontrolörler             & Arduino Uno - ESP32         & —         \\
    Bluetooth Modülü          & HC-05               & UART      \\
    Sıcaklık Sensörü          & DS18B20             & 1-Wire    \\
    Işık Sensörleri (jest)       & LDR çifti           & GPIO/ADC  \\
    Röle Modülü               & 5 V Röle            & GPIO      \\
    OLED Ekran                & 0.96" SSD1306       & I2C       \\
    \bottomrule
  \end{tabular}
\end{table}

% ════════════════════════════════════════════════════════════════════════════
\section{Özellik 1: Soğuk Algılama ve Müzik Açma}
% ════════════════════════════════════════════════════════════════════════════

\subsection{Çalışma Prensibi}

\textbf{DS18B20} sıcaklık sensörü, \textbf{1-Wire} protokolü üzerinden \textbf{Arduino Uno}'ya periyodik olarak ölçülen sıcaklık derecesini iletmektedir. Ölçüm döngüsü, ana program akışı içerisinde yazılımsal bir gecikme (\textbf{delay(2000)}) aracılığıyla 2 saniyede bir tetiklenmektedir. Sensörden okunan dijital veriler, harici bir ADC dönüşüme gerek kalmadan \textbf{DallasTemperature} kütüphanesi kullanılarak doğrudan santigrat dereceye çevrilmektedir.
\\ \\
Dönüştürülen bu değer, sistemde belirlenen hedef sıcaklık eşiği (örneğin 10°C) ile kıyaslanmaktadır. Sıcaklık limitinin altına inildiğinde ve ilgili bayrak değişkeni (\textbf{sarkiCalindi} ) kontrol edildiğinde sistem bekleme durumundan çıkarak \textbf{Bluetooth} modülünü tetikler.

\subsection{Devre Şeması}


\begin{figure}[H]
  \centering
    \resizebox{\textwidth}{!}{\includegraphics{a.png}}
  \caption{Soğuk Algılama ve Müzik Açma Devre Şeması}
  \label{fig:devre_semasi}
\end{figure}
\subsection{Bluetooth Haberleşmesi (HC-05 / UART)}

Sistemdeki haberleşme \textbf{SoftwareSerial} kütüphanesi kullanılarak \textbf{D10 (RX)} ve \textbf{D11 (TX) }pinleri üzerinden 9600 baud hızında yapılmıştır. Sıcaklık eşiği aşıldığında \textbf{HC-05} modülü üzerinden "\textbf{PLAY\_MUSIC}" komutu gönderilmektedir. Bilgisayar ortamında arka planda çalışan Python uygulaması (\textbf{pyserial} ve \textbf{spotipy} kütüphaneleri ile) bu portu dinlemekte ve komutu yakaladığı an token ile yetkilendirilmiş API üzerinden Spotify'da olan şarkıyi otonom olarak başlatmaktadır.

\subsection{Kalibrasyon ve Test Sonuçları}

Sensörden gelen veriler, Arduino IDE Seri Port Ekranı ve I2C protokolüyle çalışan \textbf{ 0.96" OLED}  ekran üzerine yansıtılıp, bu veriler kontrol edildi.. Oda sıcaklığında yapılan testlerde hedef sıcaklık limiti kod üzerinde değiştirilerek ve ek olarak sıcaklık sensörü soğuk suya koyularak,  farklı testler gerçekleştirilmiştir. Bu testler sonucunda \textbf{Bluetooth UART} komutunun terminale kayıpsız ulaştığını ve bilgisayar üzerindeki Python kodunun müziği anında başlattığı doğrulanmıştır.

% ════════════════════════════════════════════════════════════════════════════
\section{Özellik 2: Temassız Jest Kontrolü ile Müzik Değiştirme}
% ════════════════════════════════════════════════════════════════════════════

\subsection{Çalışma Prensibi}

İki adet LDR sensörü birbirinden tutarlı bir mesafede
yerleştirilmiştir. Bu sensörler, sistem başlarken ve her 30 saniyede ortam parlaklığının eşik değerini ölçme görevini üstlenir. Jest sırasında her iki LDR sensörü için o anki ölçülen parlaklık eşik değerinin altındaysa jest algılanır ve hangi LDR sensörünün ilk tetiklendiğine göre komut belirlenir:

\begin{itemize}[leftmargin=*]
  \item Sol $\to$ Sağ: Bir sonraki şarkıya geç $\to$ \texttt{NEXT}
  \item Sağ $\to$ Sol: Önceki şarkıya geç $\to$ \texttt{PREV}
\end{itemize}

\subsection{Jest Algılama Kontrolleri}

Ortamdaki ani ışık değişimlerinin yanlış jest olarak algılanmaması için bazı parametreler belirlenmiştir:

\begin{enumerate}[leftmargin=*, label=\arabic*.]
  \item İlk LDR tetiklendiğinde donanımsal zamanlayıcı başlatılır.
  \item İkinci LDR belirlenen $t_{\max} = 800\,\mathrm{ms}$ penceresi içinde
        tetiklenirse jest geçerli sayılır.
  \item Pencere süresi aşılırsa olay görmezden gelinir ve zamanlayıcı sıfırlanır.

Kodda kullanılan 800 ms test sonucu gözlemlerle belirlenmiştir.
\end{enumerate}


\subsection{Spotipy}

LDR sensörleri tarafından algılanan komut OLED ekrana basılır ve aynı zamanda çalışan python koduyla Spotify API ile iletişim başlatılır. \texttt{NEXT} komutu ile birlikte \textbf{halihazırda çalan şarkı} değiştirilir ve bir sonraki şarkıya geçilir. \texttt{PREV} komutu ise bir önceki şarkıya geçişi sağlar. LDR sensörleri tarafından algılanan komut ve Spotify API iletişimi Bluetooth bağlantısı sayesinde gerçekleşir.

\begin{figure}[H]
  \centering
    \resizebox{\textwidth}{!}{\includegraphics{b.png}}
  \caption{Temassız Jest Kontrolü ile Müzik Değiştirme Devre Şeması}
  \label{fig:devre_semasi2}
\end{figure}
% ════════════════════════════════════════════════════════════════════════════
\section{Özellik 3: Akustik Lamba Kontrolü}
% ════════════════════════════════════════════════════════════════════════════

\subsection{Çalışma Prensibi}

KY-037 mikrofon modülü analog çıkışı (A0) üzerinden Arduino Uno'ya sürekli ses verisi iletmektedir. Arduino her döngüde \texttt{analogRead()} ile 0--1023 aralığında bir değer okur ve bu değeri yazılımsal eşik (45) ile karşılaştırır. Eşiği aşan ses algılandığında alkış sayacı güncellenir; 1 saniyelik karar penceresi sonunda alkış sayısına göre işlem yapılır.

Tek alkışta lamba açılıp kapatılır. Çift alkışta ise lamba açıksa parlaklık kademesi değiştirilir. Parlaklık kontrolü Arduino'nun D9 PWM pini üzerinden IRF520 MOSFET modülüne sinyal gönderilerek sağlanır. MOSFET modülü 12V DC adaptörden beslenen LED lambayı sürer.

\subsection{Tek/Çift Alkış Algoritması}

\begin{enumerate}[leftmargin=*, label=\arabic*.]
  \item Ses eşiği aşıldığında \texttt{firstClap} zamanı kaydedilir, \texttt{clapCount} 1 yapılır.
  \item 1 saniyelik pencere içinde ikinci alkış gelirse \texttt{clapCount} 2 yapılır.
  \item 1 saniye dolunca karar verilir: 1 ise aç/kapat, 2 ise kademe değişimi.
  \item Ardışık darbeleri önlemek için 200 ms debounce uygulanır.
\end{enumerate}

\subsection{Parlaklık Kademeleri}

Eşik değeri belirlendikten sonra ölçülen değer eşik değerinin \%80'nin altında ise jest başladı/devam etmekte olarak algılanır ve LDR sensörleri tetiklenir.


% ════════════════════════════════════════════════════════════════════════════
\section{Sonuç}
% ════════════════════════════════════════════════════════════════════════════

Bu proje kapsamında, Arduino Uno ve ESP32 mikrodenetleyicileri üzerinde Bluetooth bağlantısı kullanılarak Spotify API tetikleyen bir sistem geliştirilmiştir. Üç temel işlev (sıcaklık tabanlı müzik
başlatma, temassız jest kontrolü ve akustik lamba yönetimi) birbirinden bağımsız
alt sistemler olarak tasarlanmış; ortak bir Spotify hesabı çatısı altında, her işlem başka bir sistemi tetikleyecek bir şekilde düzen kurulmuştur.



% ════════════════════════════════════════════════════════════════════════════
\appendix
\section{ Kaynak Kodu}
% ════════════════════════════════════════════════════════════════════════════

Projenin  kaynak kodları aşağıdaki GitHub deposunda yer almaktadır:\\
\url{https://github.com/sanberkkk/Micro.git}

% İsteğe bağlı: Kodları doğrudan buraya eklemek için aşağıdaki yapıyı kullanın:
% \lstinputlisting[style=cstyle, caption={main.c}]{../src/main.c}

\end{document}